%% main.tex -- arXiv-style preprint
%% Audit target: szl-holdings/lutar-lean, merged PR #189
%% Frozen commit: 8de25baf1d5adcc11238d13e17a9a7eaaf05af6d
\documentclass[10pt,letterpaper]{article}
\usepackage[margin=1in]{geometry}
\usepackage{amsmath,amssymb,amsthm}
\usepackage{booktabs,longtable,tabularx,array}
\usepackage{listings}
\usepackage{xcolor}
\usepackage{microtype}
\usepackage[colorlinks=true,linkcolor=blue,citecolor=blue,urlcolor=blue]{hyperref}
\definecolor{codegray}{RGB}{248,248,248}
\definecolor{rulegray}{RGB}{210,210,210}
\lstset{basicstyle=\ttfamily\footnotesize,breaklines=true,
  breakatwhitespace=false,frame=single,framerule=0.3pt,
  rulecolor=\color{rulegray},backgroundcolor=\color{codegray},
  columns=fullflexible,keepspaces=true,xleftmargin=0.8em,
  framexleftmargin=0.5em,showstringspaces=false}
\newtheorem{theorem}{Theorem}[section]
\newtheorem{proposition}[theorem]{Proposition}
\newtheorem{corollary}[theorem]{Corollary}
\newtheorem{lemma}[theorem]{Lemma}
\theoremstyle{definition}
\newtheorem{definition}[theorem]{Definition}
\newtheorem{remark}[theorem]{Remark}
\newtheorem{example}[theorem]{Example}
\newcolumntype{Y}{>{\raggedright\arraybackslash}X}
\newcolumntype{L}[1]{>{\raggedright\arraybackslash}p{#1}}
\newcommand{\Lam}{\Lambda}
\newcommand{\Lgraph}{\Lambda_{\mathrm{graph}}}
\newcommand{\coretrio}{\{\texttt{propext},\,\allowbreak
  \texttt{Classical.choice},\,\allowbreak\texttt{Quot.sound}\}}
\newcommand{\none}{\(\varnothing\)}
\newcommand{\decl}[1]{\nolinkurl{#1}}
\newcommand{\sha}[2]{\texttt{#1}\allowbreak\texttt{#2}}
\newcommand{\commit}{\sha{8de25baf1d5adcc11238}{d13e17a9a7eaaf05af6d}}
\newcommand{\shortcommit}{\texttt{8de25baf}}
\newcommand{\baseline}{\texttt{c7c0ba17}}
\setlength{\LTpre}{0.5em}
\setlength{\LTpost}{0.5em}
\setlength{\tabcolsep}{4pt}
\emergencystretch=2em
\title{\textbf{A Lean 4 Audit of Graph-Shaped Governance:}\\
Permutation-Safe Score Aggregation, \(1\)-WL Factoring,\\
Geometric Contraction, and Finite-Frontier Bounds}
\author{Stephen P. Lutar Jr.\thanks{SZL Holdings, Inc. ORCID
  \href{https://orcid.org/0009-0001-0110-4173}{0009-0001-0110-4173}.
  Correspondence: \texttt{stephenlutar2@gmail.com}. License CC-BY-4.0.
  Companion program DOI:
  \href{https://doi.org/10.5281/zenodo.19944926}{10.5281/zenodo.19944926}.
  Audited repository:
  \href{https://github.com/szl-holdings/lutar-lean}{github.com/szl-holdings/lutar-lean}.}}
\date{July 2026}
\begin{document}
\maketitle
\begin{abstract}
Formal verification can establish that a theorem follows from its premises
without establishing that the formal model captures the operational claim
attached to it. This distinction is acute for graph-shaped AI governance:
putting a graph field in a record does not make an aggregate topology-sensitive,
and naming a generic factoring lemma after Weisfeiler--Lehman (WL) refinement
does not construct the required factorization. We conduct a source-level and
kernel-evidence audit of a Lean~4 artifact merged as pull request \#189 at commit
\commit. The focal artifact contains 17 printed declarations across
\texttt{GraphLambda}, \texttt{GraphSubstrate}, and \texttt{MetricSpectral}; we
also reconcile eight contemporaneous information-theoretic declarations that
an earlier manuscript incorrectly described as deferred. A public CI run built
the same Git tree and printed the axiom dependencies for all 25 declarations.
The positive result is reproducible proof evidence: none of the printed
declarations depends on \texttt{sorryAx}; seven are axiom-free, four use only
\texttt{propext}, one uses \texttt{propext} and \texttt{Quot.sound}, and
thirteen use the standard Mathlib trio \coretrio. The semantic audit is more
qualified. The implemented
\(\texttt{vertexLambda}(v)=\texttt{Lutar.}\Lam(9,\texttt{scores}(v))\) never
reads adjacency, so \(\Lgraph\) is invariant under score-preserving
reindexing but is topology-blind. The graph-isomorphism witnesses contain an
\texttt{edge\_pres} field that their proofs do not consume. The declaration
\texttt{gnn\_le\_wl} proves a correct implication once a round-wise
factorization is supplied, but it neither defines WL refinement nor proves that
an \texttt{MPSystem} factors through it. The finite-frontier declarations prove
termination of a natural-number countdown, not yet of a graph traversal.
Metric declarations correctly provide coordinate non-expansion and geometric
iteration bounds, subject to explicit premises; they do not formalize a
spectral gap. We turn these findings into a reusable non-vacuity protocol,
countermodels, a declaration-by-declaration evidence ledger, deployment
guidance, and a reviewer checklist. The contribution is a precise boundary
between kernel-checked proof validity and graph-semantic adequacy.
\end{abstract}
\noindent\textbf{Keywords:} Lean 4, formal verification, graph neural networks,
Weisfeiler--Lehman, non-vacuity, proof audit, AI governance, metric contraction,
permutation invariance, reproducibility.

\section{Introduction}\label{sec:intro}

Graph language is attractive in governed AI. Services become vertices,
communication channels become edges, receipts form a directed acyclic graph,
and trust is summarized over the resulting mesh. Each metaphor suggests a
specific assurance claim: a verdict should not change when vertices are
renamed; message passing should have a known expressivity ceiling; an audit
walk should terminate; and iterative updates should remain stable. Formal
verification appears to offer an unusually strong answer because a proof
assistant checks every term against a small kernel.

There are, however, two independent questions:
\begin{enumerate}
\item \emph{Proof validity:} does the declaration elaborate, avoid
  \texttt{sorryAx}, and depend only on disclosed axioms?
\item \emph{Model adequacy:} do its definitions and hypotheses express the
  graph-semantic or operational claim used to motivate it?
\end{enumerate}
The first question is answered by the Lean kernel and a reproducible build. The
second requires an audit of definitions, consumed premises, countermodels, and
the refinement link between an abstract transition and deployed code. A theorem
may be perfectly valid while its surrounding prose exceeds its formal content.
For example, every constant function is invariant under graph isomorphism. That
fact alone does not make the function a useful graph statistic.

This paper applies both questions to the Wave-6 graph substrate in
\texttt{szl-holdings/lutar-lean}. The audited merge, pull request \#189, is
frozen at \shortcommit. Its source introduces graph-level trust aggregation,
an abstract message-passing system, a declaration named
\texttt{gnn\_le\_wl}, a natural-number frontier model, finite relabelling
lemmas, metric-coordinate bounds, and a geometric iteration theorem. An earlier
draft treated these as eleven theorems and summarized them as a unified
kernel-verified graph substrate. That count does not agree with the printed
declarations, and several summaries imported meaning from comments rather than
from types.

\subsection{Contributions}
The paper makes five contributions.
\begin{enumerate}
\item \textbf{A frozen, reproducible evidence target.} We identify the merged
  commit, the pull-request head, their identical Git tree, the exact toolchain,
  the public CI run, and the blob identifiers of every audited source file.
\item \textbf{A complete declaration ledger.} We enumerate 25 declarations:
  17 in the graph and metric scope and eight information-theoretic declarations
  needed to correct the historical record. Each entry records proof shape,
  premises consumed, axiom output, and semantic scope.
\item \textbf{A topology-sensitivity audit.} We show that
  \texttt{vertexLambda} ignores the graph field. Consequently,
  \(\Lgraph\) is a score multiset aggregate stored in a graph-shaped record,
  not a neighborhood aggregate.
\item \textbf{A claim-by-claim scope correction.} We restate the exact content
  of the WL factoring lemma, countdown termination lemmas, relabelling lemmas,
  Frechet/Kuratowski results, and geometric iteration theorem.
\item \textbf{A non-vacuity protocol.} We propose dependency,
  premise-sensitivity, discriminating-witness, negative-case, refinement, and
  deployment checks.
\end{enumerate}

\subsection{What this paper does not claim}
We do not claim a novel WL theorem, a new spectral mixing theorem, a complete
receipt-DAG verification, or topology-sensitive trust aggregation. We do not
reclassify any Wave-6 result as part of the repository's frozen locked baseline
\baseline. All audited declarations remain experimental-tier artifacts. We do
not infer deployed behavior from a successful proof build, and we report no
performance benchmark. This paper is distinct from broader work on
evidence-typed service readiness: its object is the semantic adequacy of a
specific Lean graph model and its exact proof ledger.

\section{Audit target and evidence method}\label{sec:method}

\subsection{Frozen snapshot}
The audited object is merge commit \commit, with subject
\emph{prove-wave6: graph+info substrate (F-G1..F-G6) onto Mathlib v4.18.0
(\#189)}. It was created on 6 June 2026 and has parent
\sha{9c705fbf8409d7cd370e}{de4abb04a3e931d915fb}. The successful pull-request
workflow ran on head
\sha{be376b40766a29d59549}{f7f22a113758e5c475e5}. The PR head and merge commit
have the identical Git tree
\sha{d9db68448b337949aeae}{be37409b453abad908b6}. The CI evidence therefore
covers the exact merged content even though the displayed head SHA differs
from the merge SHA.

The pinned toolchain is Lean \texttt{v4.18.0}. The manifest pins Mathlib
\texttt{v4.18.0} at
\sha{aa936c36e8484abd3005}{77139faf8e945850831a}. The public workflow
\emph{Lake build (gate + numbers)}, run
\href{https://github.com/szl-holdings/lutar-lean/actions/runs/27066757406}{27066757406},
completed successfully and emitted \texttt{\#print axioms} output for the
declarations inventoried below.

\begin{table}[ht]
\centering
\small
\begin{tabularx}{\textwidth}{@{}L{0.31\textwidth}L{0.39\textwidth}Y@{}}
\toprule
Source file & Git blob at \shortcommit & Role in this audit \\
\midrule
\nolinkurl{Lutar/GraphLambda.lean} &
\sha{b35ecea322be7f7aed7b}{78e03f3f50c824210d6c} &
Graph-shaped execution and score aggregation \\
\nolinkurl{Lutar/Wave6/GraphSubstrate.lean} &
\sha{d79168ec78b2e484433d}{3266d4fc82e222cb4275} &
Factoring, countdown, and relabelling cores \\
\nolinkurl{Lutar/Wave6/MetricSpectral.lean} &
\sha{76c74d8ea8acc7226eca}{b5a2445c400263102b9b} &
Metric coordinates and iteration bounds \\
\nolinkurl{Lutar/Wave6/InfoSubstrate.lean} &
\sha{cabf498f77a016d255c2}{3783b22fa03e8eb36657} &
Discrete information cores; context correction \\
\nolinkurl{Lutar/Wave6/SubGaussianKL.lean} &
\sha{aa8cda6102d0f5793763}{3eabc9315a3cea7d64f4} &
Three Mathlib re-exports; context correction \\
\bottomrule
\end{tabularx}
\caption{Immutable source identifiers for the audited artifact.}
\label{tab:blobs}
\end{table}

\subsection{Three evidence layers}
\begin{definition}[Kernel evidence]
A declaration has kernel evidence when the pinned Lean toolchain elaborates and
checks it, its \texttt{\#print axioms} output is captured, and the relevant
source does not introduce \texttt{sorryAx} into that declaration.
\end{definition}

\begin{definition}[Semantic evidence]
A declaration has semantic evidence for a named system claim when its
definitions mention the relevant system structure, its proof consumes premises
that encode that structure, and countermodels excluded by the prose claim are
also excluded by the formal statement.
\end{definition}

\begin{definition}[Refinement evidence]
An abstract theorem has refinement evidence for an implementation when a
checked or independently tested relation maps implementation states and steps
to the abstract model and preserves the theorem's hypotheses.
\end{definition}

Kernel evidence is neither downgraded nor treated as sufficient for the other
layers. A valid theorem about a countdown is useful; it becomes evidence for a
deployed queue only after a refinement argument shows that each queue step
decreases that countdown and cannot create unaccounted work.

\subsection{Audit procedure}
The procedure was:
\begin{enumerate}
\item detach at the merged SHA and record the Git tree and source blobs;
\item read definitions before theorem comments, because comments are not
  kernel-checked;
\item enumerate every \texttt{\#print axioms} line in the five selected files;
\item obtain the exact axiom output from the successful CI log;
\item inspect proof terms for reflexivity, direct hypothesis discharge,
  re-export, arithmetic, induction, and reindexing;
\item record which structure fields and hypotheses the proof actually consumes;
\item construct minimal countermodels for stronger prose interpretations;
\item distinguish the theorem itself from the missing refinement obligation;
\item compile and visually inspect this manuscript from its source.
\end{enumerate}
The procedure deliberately avoids theorem-count marketing. Supporting lemmas,
wrappers, direct re-exports, and substantive inductions are all valid
declarations, but they do not carry identical evidentiary weight.

\section{System and threat model}\label{sec:system}
\subsection{Four formal components}
\paragraph{Graph-shaped score record.}
\texttt{GraphExecution} packages a finite vertex type, decidable equality, a
\texttt{SimpleGraph}, a nine-axis score vector per vertex, and a bound showing
every score is at most one. A field's presence does not imply that subsequent
definitions read it.

\paragraph{Abstract message passing.}
\texttt{MPSystem} has an aggregate function
\(\texttt{aggV}:(V\to L)\to V\to A\) and an update function
\(\texttt{updV}:L\to A\to L\). The recursive \texttt{mpRun} is deterministic
by definition. Separately, \texttt{gnn\_le\_wl} is stated for arbitrary
functions \(c\) and \(f\), plus an explicit decoder \(\phi\) witnessing that
\(f\) factors through \(c\).

\paragraph{Finite countdown.}
\texttt{stepRemaining} decrements a natural number when it is positive.
\texttt{iterStep} repeats that function. This model contains no vertices,
edges, queue, frontier set, reachability relation, or DAG invariant.

\paragraph{Metric iteration.}
\texttt{MetricSpectral} works in an arbitrary pseudo-metric space. It proves
coordinate inequalities, re-exports a Kuratowski isometry, and propagates a
user-supplied one-step distance bound through an iterator.

\subsection{Threats to assurance}
\begin{table}[ht]
\centering
\small
\begin{tabularx}{\textwidth}{@{}L{0.11\textwidth}L{0.28\textwidth}Y@{}}
\toprule
ID & Threat & Required evidence \\
\midrule
T1 & Vertex relabelling changes a verdict & Permutation or isomorphism invariance \\
T2 & Topology changes while a graph-aware verdict remains unchanged &
Discriminating graph witnesses and live adjacency dependency \\
T3 & A factoring premise is presented as a universal expressivity theorem &
Definition of WL, model class, and construction of the factorization \\
T4 & An implementation can add work although the abstract measure only falls &
Simulation/refinement proof for the real transition relation \\
T5 & Convergence language hides an assumed one-step contraction &
Evidence for the metric, operator, factor, and one-step inequality \\
T6 & Counts mix reflexivity, wrappers, re-exports, and substantive results &
Declaration-level inventory and proof-shape classification \\
T7 & A branch head or current main silently replaces the reviewed snapshot &
Merged commit, tree, blobs, toolchain, and CI identifiers \\
\bottomrule
\end{tabularx}
\caption{Assurance threats addressed by the audit.}
\label{tab:threats}
\end{table}

\subsection{Adversary capabilities and scope}
For T1 and T2, an adversary may rename vertices, preserve scores while changing
edges, or choose a graph whose structure is operationally different but whose
score multiset is identical. For T3, it is enough to omit the factoring witness
while retaining the theorem's label in documentation. For T4, a scheduler may
retry, enqueue, or expand a frontier in ways absent from a natural-number
abstraction. For T5, the update map may fail the assumed contraction. T6 and T7
are provenance threats: selective counting and snapshot drift can overstate a
valid artifact without altering its source. We do not model malicious changes
to the Lean kernel, compiler, or Git object database.

\section{Graph-shaped aggregation and topology blindness}\label{sec:lambda}

\subsection{The implementation}
The decisive definitions are short. The record contains a graph:
\begin{lstlisting}[language={},caption={Selected fields of \texttt{GraphExecution}.}]
structure GraphExecution where
  V : Type
  [V_fintype : Fintype V]
  [V_dec : DecidableEq V]
  graph : SimpleGraph V
  scores : V -> Axes 9
  bounded : forall v i, scores v i <= 1
\end{lstlisting}
but the per-vertex value is:
\begin{lstlisting}[language={},caption={The exact score-only definition at the audited snapshot.}]
noncomputable def vertexLambda
    (e : GraphExecution) (v : e.V) : NNReal :=
  Lutar.Lambda 9 (e.scores v)
\end{lstlisting}
The source uses the Unicode identifier \(\Lam\); Listing~2 transliterates it as
\texttt{Lambda} for portable rendering. There is no adjacency query, neighbor
enumeration, closed-neighborhood filter, degree term, or message exchange.
The graph-level value is the geometric mean of the per-vertex score aggregates:
\[
\Lgraph(e)=
\begin{cases}
0, & |V|=0,\\[3pt]
\left(\prod_{v\in V}\Lam_9(\texttt{scores}(v))\right)^{1/|V|},
& |V|>0.
\end{cases}
\]
This is a well-defined score-multiset functional. It is permutation invariant
for the same reason any symmetric product is permutation invariant. It is not,
at this snapshot, a neighborhood aggregate.

\subsection{A source-level topology-blindness result}
\begin{proposition}[Topology blindness]\label{prop:blind}
Fix a finite carrier \(V\) and score assignment \(s:V\to\texttt{Axes}\ 9\).
Let \(e_1\) and \(e_2\) be \texttt{GraphExecution} values with that carrier and
score assignment but arbitrary simple graphs \(G_1\) and \(G_2\). Then
\(\texttt{vertexLambda}(e_1,v)=\texttt{vertexLambda}(e_2,v)\) for every \(v\),
and \(\Lgraph(e_1)=\Lgraph(e_2)\).
\end{proposition}
\begin{proof}
After unfolding \texttt{vertexLambda}, both sides are
\(\Lam_9(s(v))\). The graph field is absent. The products range over the same
finite carrier and have pointwise equal factors; the cardinality and exponent
are equal. The empty-carrier cases also agree. This is a metatheoretic
consequence of the audited definition, not a separately named declaration in
the snapshot.
\end{proof}

\begin{example}[Minimal discriminating countermodel]\label{ex:empty-path}
Take three vertices and assign the same valid nine-axis score vector to each.
Let \(G_1\) be the empty graph and \(G_2\) be a three-vertex path. They are not
isomorphic because their edge counts and degree multisets differ. Nevertheless,
Proposition~\ref{prop:blind} gives identical per-vertex and global
\(\Lam\) values. Thus \(\Lgraph\) cannot detect a total loss of graph
connectivity when scores remain fixed.
\end{example}
The countermodel does not refute the Lean declarations. It refutes the stronger
interpretation that the value aggregates closed neighborhoods or responds to
topology. This is why theorem validity and semantic adequacy must be reported
separately.

\subsection{Unused edge preservation}
The automorphism witness is stronger than the proof needs:
\begin{lstlisting}[language={},caption={The witness contains edge and score preservation.}]
structure LambdaAutomorphism (e : GraphExecution) where
  toFun : e.V -> e.V
  bij : Function.Bijective toFun
  edge_pres :
    forall v w,
      e.graph.Adj v w <->
      e.graph.Adj (toFun v) (toFun w)
  score_pres :
    forall v, e.scores v = e.scores (toFun v)
\end{lstlisting}
Yet the per-vertex proof is:
\begin{lstlisting}[language={},caption={The proof consumes score preservation, not edge preservation.}]
theorem vertexLambda_automorphism_invariant
    (phi : LambdaAutomorphism e) (v : e.V) :
    vertexLambda e (phi.toFun v) = vertexLambda e v := by
  unfold vertexLambda
  rw [<- phi.score_pres v]
\end{lstlisting}
The ASCII arrow represents Lean's leftward rewrite. The
\texttt{edge\_pres} field is not referenced. Downstream product and graph-level
proofs use score transport, the vertex equivalence, and cardinality; they still
do not require adjacency. The same issue occurs for the cross-graph
\texttt{LambdaIso} witness.

\begin{definition}[Live structure field]
A structure field \(x\) is live for a functional \(F\) if there exist two
well-typed inputs differing in \(x\), while all other semantically relevant
fields are held fixed, such that \(F\) returns different values.
\end{definition}
By Proposition~\ref{prop:blind}, \texttt{graph} is not live for
\(\Lgraph\). The proof witness can still require edge preservation, but that
requirement is phantom structure: it narrows the theorem's input class without
explaining the output. A stronger premise cannot compensate for a definition
that ignores the field.

\subsection{What the four GraphLambda declarations establish}
The four printed declarations remain correct:
\begin{enumerate}
\item \decl{Lambda_graph_automorphism_invariant} shows equality of the product
  after applying a score-preserving vertex map;
\item \decl{Lambda_graph_invariant_under_automorphism} wraps that equality in
  the nonempty graph-level expression;
\item \decl{prod_vertexLambda_iso_invariant} transports the product across a
  vertex equivalence with score preservation; and
\item \decl{Lambda_graph_iso_invariant} handles equal cardinality and the empty
  case to obtain equality of graph-level values.
\end{enumerate}
The source identifiers use the Greek \(\Lam\) where this PDF uses
\texttt{Lambda} in breakable names. All four depend on \coretrio according to
CI. Their honest collective statement is:
\begin{quote}
The geometric mean of per-vertex nine-axis score aggregates is invariant under
score-preserving reindexing. When a graph isomorphism is also supplied, the
same equality follows, although adjacency preservation is not needed by the
current functional.
\end{quote}
This is useful for serialization, sharding, and stable audit display. It should
not be labeled a topology-aware trust score.

\subsection{A topology-sensitive successor}
A future graph-sensitive definition must make adjacency live. One candidate is
a closed-neighborhood aggregate
\[
\Lam_N(e,v)=
\operatorname{GM}\{\Lam_9(\texttt{scores}(u)):
u=v\ \lor\ \texttt{Adj}_e(v,u)\},
\]
followed by a graph-level aggregation over \(v\). A rigorous Lean
implementation would need:
\begin{enumerate}
\item a finite closed-neighborhood \texttt{Finset} with decidable adjacency;
\item an explicit policy for isolated vertices and an empty carrier;
\item transport of neighborhood membership through \texttt{edge\_pres};
\item reindexing of the neighborhood product under an equivalence;
\item a discriminating witness showing some edge change can alter the value;
\item bounds compatible with the underlying \(\Lam_9\) score.
\end{enumerate}
Isomorphism invariance would then consume adjacency preservation for a reason.
The discriminating witness need not show that every edge change alters the
score; it establishes that topology can matter.

\section{The \(1\)-WL factoring declaration}\label{sec:wl}

\subsection{Exact theorem shape}
The file defines an abstract message-passing recurrence:
\begin{lstlisting}[language={},caption={The abstract message-passing system.}]
structure MPSystem (V L A : Type _) where
  aggV : (V -> L) -> V -> A
  updV : L -> A -> L

def mpRun (S : MPSystem V L A) (init : V -> L) :
    Nat -> V -> L
  | 0 => init
  | t + 1 => fun v =>
      S.updV (mpRun S init t v)
        (S.aggV (mpRun S init t) v)
\end{lstlisting}
The printed declaration \texttt{mpRun\_det} has the conclusion
\(\texttt{mpRun}\ S\ \texttt{init}\ t =
\texttt{mpRun}\ S\ \texttt{init}\ t\) and proof \texttt{rfl}. It is
axiom-free, but it is reflexivity rather than a determinism theorem comparing
two runs or transition relations. The recursion itself is deterministic by
construction; the named declaration adds no nontrivial constraint.

The declaration labeled \(\mathrm{GNN}\leq 1\text{-WL}\) is:
\begin{lstlisting}[language={},caption={The exact logical core of \texttt{gnn\_le\_wl}.}]
theorem gnn_le_wl
    (c : Nat -> V -> C) (f : Nat -> V -> L)
    (phi : Nat -> C -> L)
    (hfac : forall t v, f t v = phi t (c t v))
    (hwl : c t u = c t v) :
    f t u = f t v := by
  rw [hfac t u, hfac t v, hwl]
\end{lstlisting}
This is an axiom-free and correct factoring lemma. If \(f_t=\phi_t\circ c_t\),
equal \(c_t\)-values imply equal \(f_t\)-values. The theorem does not mention
\texttt{MPSystem}, \texttt{mpRun}, adjacency, neighborhoods, multisets, an
injective aggregator, or a definition of WL refinement. The name
\texttt{hwl} does not make the arbitrary function \(c\) a WL coloring.

\subsection{The missing factorization obligation}
The classical upper-bound argument for message-passing GNNs states, under an
appropriate model class, that representations are constant on \(1\)-WL color
classes~\cite{xu2019gin,morris2019wl}. A formal proof needs to construct the
decoder \(\phi_t\) or an equivalent relation by induction. In the audited
declaration, that construction is an input:
\[
\underbrace{\forall t,v,\quad f_t(v)=\phi_t(c_t(v))}_{\texttt{hfac}}
\quad\Longrightarrow\quad
c_t(u)=c_t(v)\Longrightarrow f_t(u)=f_t(v).
\]
The implication is useful once \texttt{hfac} is established. It is not itself
evidence that every \texttt{MPSystem} satisfies \texttt{hfac} for a formal
\(1\)-WL coloring.

\begin{example}[Why the premise matters]
Let \(c_t\) assign the same color to every vertex and let \(f_t\) assign
different labels to two vertices. Then \(c_t(u)=c_t(v)\) while
\(f_t(u)\neq f_t(v)\). No decoder \(\phi_t\) can satisfy \texttt{hfac}; the
theorem is simply inapplicable. The example does not contradict the theorem.
It shows that the factoring witness carries the entire expressivity content.
\end{example}

\subsection{What a complete formal upper bound would require}
A complete artifact could proceed in five stages.
\begin{enumerate}
\item Define finite graphs and initial colors in the theorem's namespace.
\item Define one round of \(1\)-WL using an injective encoding of the current
  color and the multiset of neighbor colors.
\item Define the admissible message-passing class: shared update functions,
  permutation-invariant neighbor aggregation, and no vertex identifiers or
  extra positional features unless included in initial colors.
\item Prove by induction that equal WL colors imply equal message-passing
  labels at each round. Equivalently, construct the round-wise decoder
  \(\phi_t\).
\item Connect the abstract recurrence to the implementation that generates
  deployed graph representations.
\end{enumerate}
The first four stages establish a model theorem. The fifth establishes a system
claim. Position-aware features can intentionally alter the initial information
and may exceed the baseline discrimination considered by ordinary
\(1\)-WL~\cite{you2019pgnn}; such features must be modeled rather than hidden.

\subsection{Correct evidence label}
The correct label for \texttt{gnn\_le\_wl} at \shortcommit is
\emph{axiom-free factoring implication}. It is a reusable final step for a WL
upper-bound proof. Calling it a complete \(1\)-WL ceiling without also citing a
formal \texttt{hfac} construction would be premise laundering: a substantive
obligation would be presented as a theorem output even though it is a theorem
input.

\section{Finite-frontier and relabelling cores}\label{sec:frontier}
\subsection{Countdown termination}
The frontier model is the natural-number program:
\begin{lstlisting}[language={},caption={The entire state transition for frontier termination.}]
def stepRemaining (remaining : Nat) : Nat :=
  match remaining with
  | 0 => 0
  | n + 1 => n

def iterStep : Nat -> Nat -> Nat
  | 0, r => r
  | k + 1, r => iterStep k (stepRemaining r)
\end{lstlisting}
Four printed declarations concern this model.
\texttt{step\_lt} proves that a positive input strictly decreases.
\texttt{iterStep\_drains} proves that \(k\) iterations reach zero whenever
\(r\leq k\). \texttt{iterStep\_empties} instantiates \(r=k=n\).
\texttt{edge\_decisions\_bounded} proves
\(n\cdot\texttt{perNode}\leq n\cdot w\) from
\(\texttt{perNode}\leq w\).
These are valid finite arithmetic results. The middle proof is a genuine
structural induction. Their honest system interpretation is conditional:
\begin{quote}
If a real generation step is simulated by consuming exactly one unit of a
non-replenishing natural-number measure, then at most the initial number of
steps is needed; if every one of \(n\) items causes at most \(w\) decisions,
then there are at most \(nw\) decisions.
\end{quote}
The condition is not proven for a graph traversal at this snapshot.

\subsection{Abstraction gap}
No declaration in this model expresses:
\begin{itemize}
\item that a frontier is a set or queue of vertices;
\item that dequeued vertices are marked visited;
\item that successors are graph neighbors;
\item that the graph is acyclic;
\item that new work cannot be enqueued indefinitely;
\item that a topological or breadth-first order is respected; or
\item that an implementation step corresponds to \texttt{stepRemaining}.
\end{itemize}
Consequently, receipt-DAG termination is a proposed interpretation, not the
type of the theorem. The countdown is still valuable as the target measure for
a future refinement proof.

\begin{example}[Implementation not covered by the countdown]
Consider a worker that removes one receipt and enqueues two new receipts after
every operation. Its queue size increases even though one item was consumed.
The implementation cannot be simulated by \texttt{stepRemaining} without
discarding the newly created work, so the Lean termination theorem does not
apply. A finite DAG plus a visited-set invariant could rule this behavior out;
neither appears in the countdown model.
\end{example}

\subsection{A refinement theorem to seek}
Let \(S\) be implementation states, \(\to\) the real step relation, and
\(\mu:S\to\mathbb{N}\) a remaining-work measure. A useful bridge would prove:
\[
s\to s' \land \neg\operatorname{done}(s)
\quad\Longrightarrow\quad
\mu(s')<\mu(s),
\]
together with finite initialization, preservation of DAG and visited-set
invariants, and equivalence between \(\mu(s)=0\) and completion. Well-founded
induction would then establish termination of the actual transition relation.

\subsection{Relabelling declarations}
The file also prints:
\begin{lstlisting}[language={},caption={The core relabelling premise and conclusion.}]
theorem adj_pred_relabel_invariant
    (adj : V -> V -> Bool) (sigma tau : V -> V)
    (_h_st : forall x, tau (sigma x) = x)
    (_h_ts : forall x, sigma (tau x) = x)
    (hpres : forall a b,
      adj a b = adj (sigma a) (sigma b))
    (a b : V) :
    adj a b = adj (sigma a) (sigma b) :=
  hpres a b
\end{lstlisting}
The inverse premises are intentionally prefixed as unused. The conclusion is a
specialization of \texttt{hpres}; it does not prove adjacency preservation from
bijectivity. This declaration is an assumption projection.

\texttt{countAdj\_relabel\_invariant} is stronger as a finite list result. It
proves by induction that filtering a list of vertex pairs by \texttt{adj} has
the same length after mapping both endpoints through a function that preserves
the predicate. It does not require the function to be bijective because it
compares the given list with its image, not counts over the full finite carrier.
To derive global graph-statistic invariance, a caller must supply complete tuple
enumeration and the appropriate bijection/reindexing argument.

\subsection{Deployment-safe interpretation}
These results can safely support a bounded loop whose implementation already
consumes fixed fuel, a cost estimate once a per-item decision bound is
enforced, and a regression test that a finite edge-pair list preserves
predicate counts under a supplied adjacency-preserving map. They should not
alone certify arbitrary DAG generation, queue liveness, clustering-coefficient
invariance, or average-path-length invariance.

\section{Metric coordinates and geometric iteration}\label{sec:metric}
\subsection{Frechet and Kuratowski declarations}
\texttt{MetricSpectral} prints three embedding-related declarations.
\begin{enumerate}
\item \decl{frechet_coord_lipschitz} instantiates the reverse triangle
  inequality
  \[
  |\operatorname{dist}(x,a)-\operatorname{dist}(y,a)|
  \leq \operatorname{dist}(x,y).
  \]
\item \decl{frechet_coord_nonexpand} takes the one-sided consequence
  \[
  \operatorname{dist}(x,a)-\operatorname{dist}(y,a)
  \leq \operatorname{dist}(x,y).
  \]
\item \decl{frechet_isometric_embedding} directly re-exports Mathlib's
  Kuratowski isometry for a separable metric space into an
  \(\ell^\infty\)-style function space.
\end{enumerate}
All three have the Mathlib core dependency set \coretrio in CI. They correctly
establish per-coordinate non-expansion and an infinite-dimensional isometric
embedding. They do not establish the random-anchor lower bound, finite
dimension, or \(O(\log n)\) distortion associated with Bourgain-style
constructions~\cite{bourgain1985,llr1995}. Nor do they prove performance of a
position-aware GNN.

\subsection{Geometric propagation}
The central iteration declaration has the following logical content:
\begin{theorem}[Audited statement of geometric contraction]\label{thm:geom}
In a pseudo-metric space, let \(P:\alpha\to\alpha\), let \(\pi\in\alpha\), and
let \(\lambda\in\mathbb{R}\). Assume \(0\leq\lambda\) and
\[
\forall y,\qquad
\operatorname{dist}(P(y),\pi)
\leq \lambda\,\operatorname{dist}(y,\pi).
\]
Then for every \(t\in\mathbb{N}\) and \(x\in\alpha\),
\[
\operatorname{dist}(P^{[t]}(x),\pi)
\leq \lambda^t\,\operatorname{dist}(x,\pi).
\]
\end{theorem}
The Lean proof is induction on \(t\) plus monotonicity of multiplication by a
nonnegative number. This is a substantive and correct propagation theorem. The
formal signature does not assume \(\lambda<1\); strict convergence follows only
when a caller additionally supplies that bound. It does not define a Markov
kernel, transition matrix, eigenvalue, stationary distribution, or spectral
gap. The module name and comments provide motivation, but no spectral object
occurs in the theorem.

The signature also does not separately assume \(P(\pi)=\pi\). In a metric space
the one-step premise at \(y=\pi\) forces distance zero and hence fixedness; in a
pseudo-metric space it forces only zero distance. The safest description is
geometric distance-to-reference propagation under a global one-step bound.

\subsection{Non-increase corollary}
\decl{contraction_nonincrease} assumes \(\lambda\leq 1\) and the same one-step
bound, then proves
\(\operatorname{dist}(P(x),\pi)\leq\operatorname{dist}(x,\pi)\).
Its type does not include \(0\leq\lambda\), although satisfiability of the
one-step premise constrains negative choices whenever positive distances
occur. This is another reason to quote the exact premises instead of
paraphrasing them as a spectral theorem.

\subsection{Evidence needed for deployment}
To use Theorem~\ref{thm:geom} for a governed iterative service, one must provide:
\begin{enumerate}
\item a state space and metric corresponding to the deployed representation;
\item a deterministic or suitably lifted transition \(P\);
\item a reference state or equivalence class \(\pi\);
\item a global or invariant-region proof of the one-step inequality;
\item a factor \(0\leq\lambda<1\) if convergence is claimed;
\item a refinement argument connecting implementation iterations to
  \texttt{iterate}.
\end{enumerate}
An empirical estimate of \(\lambda\) can motivate a conjecture but is not a
substitute for a universal premise.

\section{Declaration inventory and axiom evidence}\label{sec:inventory}
\subsection{How to read the ledger}
The symbol \none means that CI reported that the declaration does not depend on
any axioms. \(\{P\}\) abbreviates \(\{\texttt{propext}\}\);
\(\{P,Q\}\) abbreviates
\(\{\texttt{propext},\texttt{Quot.sound}\}\); and
\(\{P,C,Q\}\) abbreviates \coretrio. These sets are Lean dependency reports,
not counts of mathematical assumptions in theorem hypotheses. A declaration
may be axiom-free yet rely on a strong explicit premise.

The proof-shape column is descriptive:
\begin{itemize}
\item \emph{refl/projection}: reflexivity or direct return of a premise;
\item \emph{rewrite}: substitution using supplied equalities;
\item \emph{induction/arithmetic}: a constructed finite proof;
\item \emph{reindex/wrapper}: transport across equivalences or a corollary;
\item \emph{library re-export}: direct specialization of Mathlib.
\end{itemize}
No rank is implied. A small wrapper may be the right interface, but it must not
be counted as an independent system guarantee.

\subsection{GraphLambda and GraphSubstrate}
\begin{longtable}{@{}L{0.29\textwidth}L{0.11\textwidth}L{0.18\textwidth}L{0.34\textwidth}@{}}
\caption{Printed graph declarations at \shortcommit.}\label{tab:graph-ledger}\\
\toprule
Declaration & Axioms & Proof shape & Audited scope \\
\midrule
\endfirsthead
\toprule
Declaration & Axioms & Proof shape & Audited scope \\
\midrule
\endhead
\decl{Lambda_graph_automorphism_invariant} &
\(\{P,C,Q\}\) & rewrite/product &
Score-preserving pointwise equality over the finite product; adjacency unused. \\
\decl{Lambda_graph_invariant_under_automorphism} &
\(\{P,C,Q\}\) & wrapper &
Adds the nonempty graph-level root expression; adjacency unused. \\
\decl{prod_vertexLambda_iso_invariant} &
\(\{P,C,Q\}\) & reindex &
Transports the score product across a vertex equivalence; adjacency unused. \\
\decl{Lambda_graph_iso_invariant} &
\(\{P,C,Q\}\) & wrapper/reindex &
Handles cardinality and empty cases; proves score-aggregate equality. \\
\midrule
\decl{mpRun_det} &
\none & reflexivity &
The expression equals itself; does not compare independent runs. \\
\decl{gnn_le_wl} &
\none & rewrite &
Equal colors imply equal labels after explicit factorization is supplied. \\
\decl{step_lt} &
\(\{P\}\) & arithmetic &
Positive natural-number countdown decreases by one. \\
\decl{iterStep_drains} &
\none & induction &
\(k\) countdown steps drain \(r\) when \(r\leq k\). \\
\decl{iterStep_empties} &
\none & corollary &
Instantiates the drain theorem at \(r=k=n\). \\
\decl{edge_decisions_bounded} &
\none & arithmetic &
Multiplication preserves \(\texttt{perNode}\leq w\). \\
\decl{adj_pred_relabel_invariant} &
\none & projection &
Returns the supplied pointwise adjacency-preservation premise. \\
\decl{countAdj_relabel_invariant} &
\(\{P\}\) & list induction &
Filter length is preserved on a pair list and its mapped image. \\
\bottomrule
\end{longtable}
There are 12 printed graph declarations, not eleven. Of these,
\decl{mpRun_det} is reflexive and
\decl{adj_pred_relabel_invariant} is a direct premise projection.
\decl{gnn_le_wl} is a conditional factoring implication. The table avoids
subtracting these declarations from the artifact; it prevents the count from
being mistaken for 12 independent graph guarantees.

\subsection{MetricSpectral}
\begin{table}[ht]
\centering
\small
\begin{tabularx}{\textwidth}{@{}L{0.34\textwidth}L{0.13\textwidth}L{0.18\textwidth}Y@{}}
\toprule
Declaration & Axioms & Proof shape & Audited scope \\
\midrule
\decl{frechet_coord_lipschitz} &
\(\{P,C,Q\}\) & library term &
Reverse triangle inequality for one anchor coordinate. \\
\decl{frechet_coord_nonexpand} &
\(\{P,C,Q\}\) & library corollary &
One-sided coordinate inequality. \\
\decl{frechet_isometric_embedding} &
\(\{P,C,Q\}\) & library re-export &
Kuratowski isometry for separable metric spaces. \\
\decl{geometric_contraction} &
\(\{P,C,Q\}\) & induction &
Iterates a supplied one-step distance bound for \(\lambda\geq0\). \\
\decl{contraction_nonincrease} &
\(\{P,C,Q\}\) & inequality &
One-step non-increase from the supplied bound and \(\lambda\leq1\). \\
\bottomrule
\end{tabularx}
\caption{The five printed metric declarations.}
\label{tab:metric-ledger}
\end{table}
Together, Tables~\ref{tab:graph-ledger} and \ref{tab:metric-ledger} contain the
17 declarations in this paper's focal graph-and-metric scope.

\subsection{Correcting the information-theory status}
The earlier manuscript said that Hoeffding, Azuma, and nonnegativity of
Kullback--Leibler divergence were deferred because the pinned Mathlib lacked
the required modules. That description corresponds to an older toolchain, not
the audited merge. Pull request \#189 moved to Mathlib \texttt{v4.18.0}.
At \shortcommit, \texttt{InfoSubstrate.lean} contains five discrete cores and
\texttt{SubGaussianKL.lean} contains three Mathlib-dependent analytic
re-exports. The public CI log builds both files and prints the dependencies in
Table~\ref{tab:info-ledger}.

\begin{longtable}{@{}L{0.34\textwidth}L{0.13\textwidth}L{0.18\textwidth}L{0.27\textwidth}@{}}
\caption{Information declarations included to repair the artifact record.}
\label{tab:info-ledger}\\
\toprule
Declaration & Axioms & Proof shape & Audited scope \\
\midrule
\endfirsthead
\toprule
Declaration & Axioms & Proof shape & Audited scope \\
\midrule
\endhead
\decl{dpi_postprocess_collapse} &
\none & rewrite &
An observable preserves an already supplied equality after processing. \\
\decl{dpi_count_pullback} &
\(\{P\}\) & list induction &
Mapping then filtering equals filtering by the pulled-back predicate. \\
\decl{fano_collision_forces_error} &
\(\{P,C,Q\}\) & case split &
Two distinct inputs collapsed by a decoder cannot both decode correctly. \\
\decl{coverage_conservation} &
\(\{P,Q\}\) & list induction &
Boolean covered and miscovered counts partition a finite list. \\
\decl{miscoverage_le_total} &
\(\{P\}\) & library re-export &
A filtered list is no longer than its source. \\
\decl{c3_hoeffding_sum_ge_le} &
\(\{P,C,Q\}\) & Mathlib re-export &
Specializes the independent sub-Gaussian tail theorem. \\
\decl{c4_azuma_hoeffding_sum_ge_le} &
\(\{P,C,Q\}\) & Mathlib re-export &
Specializes the conditional sub-Gaussian martingale tail theorem. \\
\decl{c5_gibbs_kl_nonneg} &
\(\{P,C,Q\}\) & Mathlib re-export &
Specializes the integral log-likelihood-ratio nonnegativity result. \\
\bottomrule
\end{longtable}

These eight declarations are not graph results and are not used to inflate the
graph contribution. They are included because an audit must correct negative
claims as carefully as positive ones. The analytic declarations were no longer
deferred at the merged snapshot.

\subsection{Axiom summary}
Across all 25 printed declarations:
\begin{center}
\begin{tabular}{@{}lr@{}}
\toprule
CI dependency output & Number of declarations \\
\midrule
\none & 7 \\
\(\{\texttt{propext}\}\) & 4 \\
\(\{\texttt{propext},\texttt{Quot.sound}\}\) & 1 \\
\coretrio & 13 \\
\midrule
Total & 25 \\
\bottomrule
\end{tabular}
\end{center}
None of these outputs contains \texttt{sorryAx} or a declared Lutar-specific
axiom. This is strong proof-validity evidence. It does not make theorem
hypotheses disappear, and it does not show that all record fields are live.

\section{Non-vacuity analysis}\label{sec:nonvacuity}
\subsection{Why axiom freedom is not enough}
An axiom-free theorem can be reflexive, a direct projection of a hypothesis, or
a conditional with an unestablished premise. This is not a flaw in Lean.
Proof assistants correctly distinguish assumptions in theorem types from
trusted axioms in the environment. Assurance reports must do the same.

We call a system interpretation \emph{non-vacuous} when the formal result rules
out at least one behavior forbidden by that interpretation and when the
relevant structure influences either the definition or a premise whose
satisfaction is independently established.

\begin{definition}[Discriminating witness]
For a proposed graph-sensitive functional \(F\), a discriminating witness is a
pair of admissible inputs that differ in graph structure and for which \(F\)
differs. It proves that topology is live, not that \(F\) is universally
sensitive.
\end{definition}

\begin{definition}[Premise closure]
A theorem application has premise closure when every premise that carries the
substantive system claim is established by a prior proof, validated contract,
or explicitly scoped measurement rather than assumed at the reporting
boundary.
\end{definition}

\subsection{Six checks}
\paragraph{1. Definition-dependency check.}
Search the definitional dependency cone for the claimed structure. If a
graph-level output never reads adjacency, graph sensitivity is absent
regardless of the theorem name.

\paragraph{2. Premise-consumption check.}
List structure fields and hypotheses referenced by the proof. An unused
\texttt{edge\_pres}, inverse witness, fixed-point comment, or DAG assumption is
evidence of a scope mismatch or a stronger-than-needed interface.

\paragraph{3. Discriminating-witness check.}
Construct the smallest positive and negative examples. For invariants, include
both a relabelling pair that must agree and a non-isomorphic pair that should be
distinguishable if topology awareness is claimed.

\paragraph{4. Assumption-ablation check.}
Identify which conclusion fails when a substantive premise is removed. For
\texttt{gnn\_le\_wl}, ablate \texttt{hfac}; for geometric contraction, ablate
the global one-step bound; for countdown termination, allow replenishment.
If the missing premise is exactly the advertised system claim, it must be
reported as an obligation, not an output.

\paragraph{5. Refinement check.}
State a simulation relation between concrete and abstract states and prove step
preservation. A theorem about \(\mathbb{N}\) does not automatically transfer to
a distributed queue, and a theorem about a pure function does not automatically
cover retries or nondeterminism.

\paragraph{6. Evidence-tier check.}
Report kernel evidence, semantic evidence, refinement evidence, and deployed
measurement separately. Never promote an experimental theorem into a frozen
baseline merely because the build is green.

\subsection{Application matrix}
\begin{table}[ht]
\centering
\scriptsize
\begin{tabularx}{\textwidth}{@{}L{0.23\textwidth}cccccY@{}}
\toprule
Claim family & Def. & Premise & Witness & Refine & Kernel & Result \\
\midrule
Score reindex invariance & yes & closed & relabel pair & n/a & yes &
Supported for the pure score aggregate. \\
Topology-sensitive trust & no edges & unused edge field & empty/path refutes & absent & yes &
Not supported at snapshot. \\
\(1\)-WL ceiling & no WL & factorization assumed & ablation refutes & absent & yes &
Factoring implication only. \\
Frontier countdown & natural no. & closed & replenish separates & absent & yes &
Abstract countdown supported. \\
Receipt-DAG termination & no graph & simulation missing & enqueue separates & absent & yes &
Not supported end to end. \\
Coordinate non-expansion & metric live & closed & metric cases & caller & yes &
Supported as a mathematical core. \\
Geometric propagation & iterate live & step bound assumed & bad map separates & caller & yes &
Supported conditionally. \\
Spectral mixing & no spectrum & gap absent & operators separate & absent & yes &
Not formalized. \\
\bottomrule
\end{tabularx}
\caption{Non-vacuity checks applied to the principal claims.}
\label{tab:nonvacuity}
\end{table}

\subsection{Negative tests as first-class artifacts}
Formal projects commonly test that intended theorems compile. They should also
test that overbroad variants fail or that counterexamples evaluate as expected.
For this substrate, useful negative artifacts include:
\begin{itemize}
\item two non-isomorphic finite graphs with identical scores and a check that
  the present \(\Lgraph\) agrees;
\item a candidate neighborhood aggregate that differs on at least one such
  pair;
\item an arbitrary coloring and nonfactoring label function demonstrating why
  \texttt{hfac} is necessary;
\item a replenishing transition that cannot satisfy a strictly decreasing
  remaining-work measure; and
\item a map that violates the one-step contraction premise.
\end{itemize}
Such tests identify the exact boundary that makes a proven theorem meaningful.

\section{Deployment implications}\label{sec:deploy}
\subsection{Names exposed to operators}
User interfaces, model cards, and service APIs should expose the narrowest
truthful names:
\begin{center}
\begin{tabularx}{0.93\textwidth}{@{}L{0.42\textwidth}Y@{}}
\toprule
Avoid & Prefer at the audited snapshot \\
\midrule
Topology-aware graph trust & Permutation-safe aggregate of vertex score vectors \\
Kernel-proven \(1\)-WL ceiling & Axiom-free factoring lemma; factorization pending \\
Receipt-DAG termination proven & Countdown termination; refinement pending \\
Spectral convergence proven & Geometric bound under a contraction premise \\
All Wave-6 results locked & Experimental declarations; baseline unchanged \\
\bottomrule
\end{tabularx}
\end{center}

\subsection{Receipt and provenance schema}
Each deployed decision receipt that relies on these artifacts should record:
\begin{itemize}
\item the theorem declaration and exact source commit;
\item the model adapter or refinement version;
\item all runtime parameters corresponding to theorem premises;
\item whether each premise was proved, configured, or measured;
\item a result tier such as kernel, semantic, refined, or observed;
\item a fail-closed outcome when a required premise lacks evidence.
\end{itemize}
For Theorem~\ref{thm:geom}, a receipt should identify the metric, \(\pi\),
\(\lambda\), invariant region, and evidence for \texttt{hcontr}. For
\texttt{gnn\_le\_wl}, it should identify the formal WL coloring and the theorem
that constructs \(\phi_t\). For countdown termination, it should identify the
concrete measure and step-simulation theorem.

\subsection{Fail-closed behavior}
Missing refinement evidence should not crash a service or silently become a
green claim. A practical policy is:
\[
\texttt{proof\_valid}
\land \texttt{premises\_closed}
\land \texttt{refinement\_current}
\quad\Longrightarrow\quad
\texttt{claim\_admissible}.
\]
If the score aggregate is requested as a topology-aware signal, the current
adapter should deny that label while still returning the score-only aggregate
under its honest name. If a queue implementation cannot show a decreasing
measure, it may run with operational monitoring but must not emit a formal
termination certificate.

\subsection{Prioritized engineering path}
\begin{enumerate}
\item Rename the current value in user-facing surfaces to make score-only
  semantics explicit.
\item Add the topology-blind counterexample as an executable regression test.
\item Implement a neighborhood-sensitive successor and prove an
  edge-consuming isomorphism theorem.
\item Connect \texttt{mpRun} to a formal WL definition and construct the
  factorization by induction.
\item Replace the countdown story with a finite traversal state, visited
  invariant, and refinement map from the production worker.
\item Define a concrete operator and metric before attaching spectral or
  mixing labels to the geometric theorem.
\item Generate a machine-readable declaration ledger from CI.
\end{enumerate}

\section{Reproducibility}\label{sec:repro}
\subsection{Source and build}
The following commands reproduce the source boundary. They require Git,
Lean \texttt{v4.18.0}, Lake, and the pinned dependencies.
\begin{lstlisting}[language={},caption={Reproduce the frozen source and build.}]
git clone https://github.com/szl-holdings/lutar-lean.git
cd lutar-lean
git switch --detach \
  8de25baf1d5adcc11238d13e17a9a7eaaf05af6d
git rev-parse HEAD
git rev-parse HEAD^{tree}
cat lean-toolchain
lake build
\end{lstlisting}
Expected identifiers are \commit and tree
\sha{d9db68448b337949aeae}{be37409b453abad908b6}. The source files contain
\texttt{\#print axioms} commands. A focused recheck can run:
\begin{lstlisting}[language={},caption={Focused module checks.}]
lake env lean Lutar/GraphLambda.lean
lake env lean Lutar/Wave6/GraphSubstrate.lean
lake env lean Lutar/Wave6/MetricSpectral.lean
lake env lean Lutar/Wave6/InfoSubstrate.lean
lake env lean Lutar/Wave6/SubGaussianKL.lean
\end{lstlisting}
Run 27066757406 reports success on head
\texttt{be376b40766a29d59549f7f22a113758e5c475e5}; its tree equals the merge
tree. The log records one linter warning in \texttt{GraphLambda.lean} for an
unused variable \texttt{h} in \texttt{Lambda\_graph\_def}. That warning does
not affect kernel validity but is part of the build record.

\subsection{Manual source checks}
A reviewer can confirm the central findings without trusting this paper:
\begin{enumerate}
\item inspect \texttt{vertexLambda} and observe that it reads
  \texttt{e.scores v} but not \texttt{e.graph};
\item inspect \texttt{vertexLambda\_automorphism\_invariant} and observe that
  its proof uses \texttt{score\_pres} only;
\item inspect \texttt{gnn\_le\_wl} and observe that \texttt{c}, \texttt{f},
  and \(\phi\) are arbitrary and \texttt{hfac} is an input;
\item inspect \texttt{stepRemaining} and observe that its state is only a
  natural number;
\item inspect \texttt{geometric\_contraction} and observe
  \texttt{hlam0} and \texttt{hcontr}, but no spectral object;
\item compare \texttt{InfoSubstrate.lean} and
  \texttt{SubGaussianKL.lean} with the old deferred claim.
\end{enumerate}

\subsection{Manuscript build}
This paper is a self-contained \texttt{article} source. The audited command is:
\begin{lstlisting}[language={},caption={Manuscript compilation used for page verification.}]
tectonic --outdir tmp/pdfs/graph-substrate \
  thesis/arxiv/graph-substrate/main.tex
\end{lstlisting}
Page count must be taken from the compiled PDF, not copied from metadata.
Visual verification should render every page and inspect the title, tables,
listings, references, page numbers, and final-page termination.

\section{Related work}\label{sec:related}
\subsection{GNN expressivity and position}
The relationship between message-passing GNNs and \(1\)-WL is established in
the graph-learning literature. Xu et al.\ characterize conditions under which
aggregation architectures match WL discrimination~\cite{xu2019gin}.
Morris et al.\ give a parallel perspective~\cite{morris2019wl}. These results
motivate the intended scope of \texttt{gnn\_le\_wl}, but a citation does not
instantiate their assumptions in Lean. The audited theorem supplies the final
factoring implication; formalizing the model class remains local work.

Position-aware GNNs use distances to anchor sets to augment local message
passing~\cite{you2019pgnn}. The coordinate inequality in
\texttt{MetricSpectral} is consistent with that motivation, while remaining
far narrower than an end-to-end P-GNN theorem. GraphGym studies architecture
design empirically~\cite{you2020graphgym}, and graph2nn relates neural
architecture topology to performance~\cite{you2020graph2nn}. Those programs
reinforce why topology should be a live formal dependency when a score is
called graph-aware.

\subsection{Metric embeddings and contraction}
The reverse triangle inequality behind each Frechet coordinate is classical.
Bourgain's theorem and its algorithmic development provide low-distortion
finite embeddings~\cite{bourgain1985,llr1995}; the audited artifact proves
neither the random-anchor lower bound nor the asymptotic statement. Its
Kuratowski re-export supplies an isometry into a large function space.

Geometric iteration bounds follow from a one-step contraction. Markov-chain
mixing theory connects spectral information to such bounds under additional
structure~\cite{levinperes2017,diaconisstroock1991}. The audited theorem starts
after the one-step inequality is supplied and does not derive it from
reversibility, eigenvalues, conductance, or a spectral gap.

\subsection{Formal verification and trusted bases}
Lean's kernel checks elaborated proof terms, while Mathlib provides reusable
mathematics~\cite{lean4,mathlib}. Projects such as seL4 and CompCert illustrate
the value of an explicit refinement chain from specification to
implementation~\cite{sel4,compcert}. Their lesson here is that theorem,
abstraction boundary, and implementation relation must be stated separately.

\section{Limitations and future work}\label{sec:limits}
\subsection{Limitations of this audit}
The audit is frozen to one merge. Later commits may strengthen definitions or
add refinement theorems; those changes require a new evidence record rather
than retroactive reinterpretation. We inspected the selected declaration
types, proofs, source dependencies, and public CI output, but did not audit the
entire transitive Mathlib implementation or Lean kernel. The axiom counts are
those printed by Lean, not a philosophical claim that classical choice or
quotients are unnecessary.

The topology-blindness proposition is derived by inspection and mathematics;
it is not checked into the audited snapshot. Adding it as a named Lean theorem
and executable finite counterexample would improve the artifact. Similarly,
the non-vacuity protocol is a review method, not a mechanized metatheory.

We do not evaluate numerical conditioning of geometric means, calibration of
nine-axis scores, fairness of trust features, or security of surrounding
services. Those are essential deployment questions but independent of whether
adjacency is read by the formal definition.

\subsection{Proof obligations for the next artifact}
\begin{enumerate}
\item \textbf{Topology liveness:} a named pair of finite graphs and scores for
  which a new neighborhood aggregate differs.
\item \textbf{Edge-consuming invariance:} a proof whose transport of
  neighborhood membership references \texttt{edge\_pres}.
\item \textbf{WL construction:} definitions of WL refinement and the allowed
  message-passing class, followed by an inductive decoder construction.
\item \textbf{Traversal refinement:} a finite graph state with frontier,
  visited set, ordering invariant, and a decreasing measure.
\item \textbf{Operator-specific contraction:} a concrete transition and metric
  with a proved one-step factor below one.
\item \textbf{Machine-readable evidence:} CI output recording declaration,
  type hash, axioms, maturity, source blob, and refinement theorem.
\end{enumerate}

\subsection{A note on the locked baseline}
The source comments distinguish the experimental Wave-6 scope from a frozen
baseline \baseline, often accompanied by the historical tuple
\(749/14/163\). This paper does not infer the semantics of that tuple or add any
audited declaration to it. A successful experimental build is evidence for the
experimental artifact only. Promotion should require a separately reviewed
release decision and the semantic/refinement checks applicable to the claim.

\section{Conclusion}\label{sec:conclusion}
The audited Lean artifact has real strengths: a frozen merge, a successful
public build, explicit axiom output, no \texttt{sorryAx} among 25 printed
declarations, correct score-reindexing equalities, finite inductions, metric
inequalities, and reusable Mathlib specializations. Those strengths become
more credible when their boundaries are stated exactly.

The central graph aggregate is topology-blind because its per-vertex definition
reads scores but never adjacency. Its isomorphism witnesses carry an unused
edge-preservation field. The declaration named after the \(1\)-WL ceiling is a
correct implication whose substantive factorization is assumed. The frontier
model is a countdown, and the geometric theorem propagates a contraction
premise without deriving a spectral gap. None of these observations invalidates
the Lean proofs. They determine what those proofs may honestly certify.

For governed AI, the practical lesson is simple: pair kernel evidence with
definition dependency, countermodels, premise closure, and implementation
refinement. Graph-shaped should not be confused with graph-sensitive, and
axiom-free should not be confused with assumption-free. A small formal core can
carry substantial assurance once the bridge from model to operational claim is
explicit.

\appendix
\section{Selected exact Lean excerpts}\label{app:lean}
This appendix preserves source fragments in a portable ASCII transliteration.
It is not a replacement for the immutable blobs in Table~\ref{tab:blobs}.

\subsection{Graph-level aggregation}
\begin{lstlisting}[language={}]
noncomputable def vertexLambda
    (e : GraphExecution) (v : e.V) : NNReal :=
  Lutar.Lambda 9 (e.scores v)

noncomputable def Lambda_graph (e : GraphExecution) : NNReal :=
  if h : Fintype.card e.V = 0 then 0
  else
    let n := Fintype.card e.V
    let prod : NNReal :=
      (Finset.univ : Finset e.V).prod (vertexLambda e)
    prod ^ ((1 : Real) / (n : Real))
\end{lstlisting}
Audit observation: the \texttt{graph} field does not occur in either body.

\subsection{Cross-graph witness and score transport}
\begin{lstlisting}[language={}]
structure LambdaIso (e1 e2 : GraphExecution) where
  toEquiv : e1.V Equiv e2.V
  edge_pres :
    forall v w,
      e1.graph.Adj v w <->
      e2.graph.Adj (toEquiv v) (toEquiv w)
  score_pres :
    forall v, e1.scores v = e2.scores (toEquiv v)

theorem vertexLambda_iso_transport
    (psi : LambdaIso e1 e2) (v : e1.V) :
    vertexLambda e1 v =
      vertexLambda e2 (psi.toEquiv v) := by
  unfold vertexLambda
  rw [psi.score_pres v]
\end{lstlisting}
Audit observation: the transport proof reads \texttt{score\_pres}; it does not
read \texttt{edge\_pres}.

\subsection{Product transport}
\begin{lstlisting}[language={}]
theorem prod_vertexLambda_iso_invariant
    (psi : LambdaIso e1 e2) :
    (Finset.univ : Finset e1.V).prod (vertexLambda e1)
      =
    (Finset.univ : Finset e2.V).prod (vertexLambda e2) := by
  exact Fintype.prod_equiv
    psi.toEquiv
    (vertexLambda e1)
    (vertexLambda e2)
    (fun v => vertexLambda_iso_transport psi v)
\end{lstlisting}
Audit observation: vertex equivalence and score transport justify reindexing.
Adjacency still does not occur.

\subsection{Factoring implication}
\begin{lstlisting}[language={}]
theorem gnn_le_wl
    {V : Type _}
    (c : Nat -> V -> C)
    (f : Nat -> V -> L)
    (phi : Nat -> C -> L)
    (hfac : forall t v, f t v = phi t (c t v))
    {t : Nat} {u v : V}
    (hwl : c t u = c t v) :
    f t u = f t v := by
  rw [hfac t u, hfac t v, hwl]
\end{lstlisting}
Audit observation: neither WL refinement nor \texttt{MPSystem} occurs in the
type. The theorem is fully general and conditional on \texttt{hfac}.

\subsection{Countdown induction}
\begin{lstlisting}[language={}]
theorem iterStep_drains :
    forall k r : Nat, r <= k -> iterStep k r = 0
  | 0, 0, _ => rfl
  | k + 1, 0, _ => by
      show iterStep k (stepRemaining 0) = 0
      have hz : stepRemaining 0 = 0 := rfl
      rw [hz]
      exact iterStep_drains k 0 (Nat.zero_le _)
  | k + 1, r + 1, h => by
      show iterStep k (stepRemaining (r + 1)) = 0
      have hs : stepRemaining (r + 1) = r := rfl
      rw [hs]
      exact iterStep_drains k r
        (Nat.le_of_succ_le_succ h)
\end{lstlisting}
Audit observation: this is a genuine induction over natural-number fuel. It
contains no graph object.

\subsection{Geometric iteration}
\begin{lstlisting}[language={},basicstyle=\ttfamily\scriptsize]
theorem geometric_contraction
    [PseudoMetricSpace alpha]
    (P : alpha -> alpha) (pi : alpha)
    (lam : Real) (hlam0 : 0 <= lam)
    (hcontr : forall y,
      dist (P y) pi <= lam * dist y pi) :
    forall (t : Nat) (x : alpha),
      dist (iterate P t x) pi
        <= lam ^ t * dist x pi := by
  intro t
  induction t with
  | zero =>
      intro x
      simp [iterate]
  | succ k ih =>
      intro x
      have h1 :
          dist (iterate P (k + 1) x) pi
            <= lam * dist (iterate P k x) pi := by
        simpa [iterate] using hcontr (iterate P k x)
      have h2 :
          lam * dist (iterate P k x) pi
            <= lam * (lam ^ k * dist x pi) :=
        mul_le_mul_of_nonneg_left (ih x) hlam0
      calc
        dist (iterate P (k + 1) x) pi
            <= lam * dist (iterate P k x) pi := h1
        _ <= lam * (lam ^ k * dist x pi) := h2
        _ = lam ^ (k + 1) * dist x pi := by ring
\end{lstlisting}
Audit observation: the induction is substantive. A global one-step bound and
nonnegative factor are premises. A factor below one and a spectral derivation
are not in the declaration.

\section{Reviewer checklist}\label{app:checklist}
\subsection{Snapshot and provenance}
\begin{itemize}
\item[$\square$] Is the reviewed object a merge commit rather than only a branch
  name?
\item[$\square$] Does CI cover the same Git tree as the merge?
\item[$\square$] Are toolchain, dependency revision, source blobs, and run URL
  recorded?
\item[$\square$] Can a reviewer reproduce each \texttt{\#print axioms} output?
\item[$\square$] Are linter warnings and build caveats retained?
\end{itemize}

\subsection{Definitions and proof dependencies}
\begin{itemize}
\item[$\square$] Does every claimed graph-sensitive value read adjacency?
\item[$\square$] Are every structure field and theorem premise used or explained?
\item[$\square$] Are reflexivity and direct hypothesis projections identified?
\item[$\square$] Are library re-exports distinguished from new formal proofs?
\item[$\square$] Are theorem hypotheses quoted alongside conclusions?
\end{itemize}

\subsection{Countermodels and non-vacuity}
\begin{itemize}
\item[$\square$] Is there a relabelling example that must preserve the output?
\item[$\square$] Is there a non-isomorphic example testing topology liveness?
\item[$\square$] Has each substantive premise been ablated in a negative test?
\item[$\square$] Does the formal property rule out a forbidden behavior?
\item[$\square$] Are counterexamples encoded or tracked as future formal work?
\end{itemize}

\subsection{Refinement and deployment}
\begin{itemize}
\item[$\square$] Is there a relation from concrete states to abstract states?
\item[$\square$] Does every implementation step preserve the relation and
  theorem premises?
\item[$\square$] Are retries, concurrency, queue growth, and nondeterminism
  represented or explicitly excluded?
\item[$\square$] Does each receipt identify theorem, commit, adapter, and
  premise evidence?
\item[$\square$] Does missing evidence fail closed at the claim boundary?
\end{itemize}

\subsection{Maturity and communication}
\begin{itemize}
\item[$\square$] Are experimental declarations separate from a locked baseline?
\item[$\square$] Is the theorem count accompanied by a declaration ledger?
\item[$\square$] Do UI labels use the narrowest truthful semantics?
\item[$\square$] Are conjectures, conditional results, and measurements
  visibly distinct?
\item[$\square$] Is page count verified from the final PDF?
\end{itemize}

\section{Claim-to-artifact crosswalk}\label{app:crosswalk}
\begin{longtable}{@{}L{0.25\textwidth}L{0.31\textwidth}L{0.36\textwidth}@{}}
\caption{Crosswalk from manuscript claims to reproducible evidence.}\\
\toprule
Manuscript claim & Primary artifact & Reviewer action \\
\midrule
\endfirsthead
\toprule
Manuscript claim & Primary artifact & Reviewer action \\
\midrule
\endhead
\(\texttt{vertexLambda}\) is score-only &
\texttt{GraphLambda.lean}, definition near the beginning of the module &
Unfold the definition; search its body for \texttt{graph} or \texttt{Adj}. \\
\(\Lgraph\) is topology-blind &
Same definition plus \texttt{Lambda\_graph} &
Instantiate fixed scores on two distinct simple graphs. \\
\texttt{edge\_pres} is unused in score transport &
\decl{vertexLambda_automorphism_invariant} and
\decl{vertexLambda_iso_transport} &
Inspect the proof bodies and their transitive calls. \\
The WL result assumes factorization &
\texttt{GraphSubstrate.lean}, \texttt{gnn\_le\_wl} &
Read \texttt{hfac}; search the theorem type for \texttt{MPSystem} and WL. \\
\texttt{mpRun\_det} is reflexive &
\texttt{GraphSubstrate.lean}, \texttt{mpRun\_det} &
Observe conclusion \(x=x\) and proof \texttt{rfl}. \\
Frontier termination is a countdown &
\texttt{stepRemaining}, \texttt{iterStep}, and drain declarations &
Search their types for a graph, queue, or transition relation. \\
The pair-count lemma is finite-list transport &
\decl{countAdj_relabel_invariant} &
Inspect the mapped list and pointwise preservation premise. \\
Metric coordinate claims are narrow &
\texttt{MetricSpectral.lean}, first three printed declarations &
Compare exact types with a finite-dimensional distortion claim. \\
Geometric iteration assumes contraction &
\texttt{geometric\_contraction} &
Read \texttt{hlam0} and \texttt{hcontr}; search for spectral data. \\
Information targets were present &
\texttt{InfoSubstrate.lean}, \texttt{SubGaussianKL.lean}, CI run &
Confirm eight printed declarations and their build output. \\
Axiom distribution is \(7/4/1/13\) &
CI run 27066757406 &
Collect the 25 relevant \texttt{\#print axioms} messages and tally sets. \\
The PR build covers merged content &
Git commits \texttt{be376b4} and \shortcommit &
Compare tree hashes; both are
\sha{d9db68448b337949aeae}{be37409b453abad908b6}. \\
\bottomrule
\end{longtable}

\begin{thebibliography}{99}
\bibitem{xu2019gin}
K.\ Xu, W.\ Hu, J.\ Leskovec, and S.\ Jegelka.
\emph{How Powerful are Graph Neural Networks?}
ICLR, 2019. arXiv:1810.00826.
\url{https://arxiv.org/abs/1810.00826}

\bibitem{morris2019wl}
C.\ Morris, M.\ Ritzert, M.\ Fey, W.\ Hamilton, J.\ Lenssen,
G.\ Rattan, and M.\ Grohe.
\emph{Weisfeiler and Leman Go Machine Learning: The Story So Far.}
AAAI, 2019. arXiv:1810.02244.
\url{https://arxiv.org/abs/1810.02244}

\bibitem{you2019pgnn}
J.\ You, R.\ Ying, and J.\ Leskovec.
\emph{Position-aware Graph Neural Networks.}
ICML, 2019. arXiv:1906.04817.
\url{https://arxiv.org/abs/1906.04817}

\bibitem{you2020graphgym}
J.\ You, R.\ Ying, X.\ Ren, W.\ Hamilton, and J.\ Leskovec.
\emph{Design Space for Graph Neural Networks.}
NeurIPS, 2020. arXiv:2011.08843.
\url{https://arxiv.org/abs/2011.08843}

\bibitem{you2020graph2nn}
J.\ You, J.\ Leskovec, K.\ He, and S.\ Xie.
\emph{Graph Structure of Neural Networks.}
ICML, 2020. arXiv:2007.06559.
\url{https://arxiv.org/abs/2007.06559}

\bibitem{bourgain1985}
J.\ Bourgain.
\emph{On Lipschitz Embedding of Finite Metric Spaces in Hilbert Space.}
Israel Journal of Mathematics, 52:46--52, 1985.
\href{https://doi.org/10.1007/BF02776078}{doi:10.1007/BF02776078}.

\bibitem{llr1995}
N.\ Linial, E.\ London, and Y.\ Rabinovich.
\emph{The Geometry of Graphs and Some of Its Algorithmic Applications.}
Combinatorica, 15:215--245, 1995.
\href{https://doi.org/10.1007/BF01200757}{doi:10.1007/BF01200757}.

\bibitem{levinperes2017}
D.\ A.\ Levin and Y.\ Peres.
\emph{Markov Chains and Mixing Times}, second edition.
American Mathematical Society, 2017.
\href{https://doi.org/10.1090/mbk/107}{doi:10.1090/mbk/107}.

\bibitem{diaconisstroock1991}
P.\ Diaconis and D.\ Stroock.
\emph{Geometric Bounds for Eigenvalues of Markov Chains.}
Annals of Applied Probability, 1(1):36--61, 1991.
\href{https://doi.org/10.1214/aoap/1177005980}{doi:10.1214/aoap/1177005980}.

\bibitem{lean4}
L.\ de Moura and S.\ Ullrich.
\emph{The Lean 4 Theorem Prover and Programming Language.}
CADE-28, 2021.
\href{https://doi.org/10.1007/978-3-030-79876-5_37}{doi:10.1007/978-3-030-79876-5\_37}.

\bibitem{mathlib}
The Mathlib Community.
\emph{The Lean Mathematical Library.}
CPP, 2020.
\href{https://doi.org/10.1145/3372885.3373824}{doi:10.1145/3372885.3373824}.

\bibitem{sel4}
G.\ Klein et al.
\emph{seL4: Formal Verification of an OS Kernel.}
SOSP, 2009.
\href{https://doi.org/10.1145/1629575.1629596}{doi:10.1145/1629575.1629596}.

\bibitem{compcert}
X.\ Leroy.
\emph{Formal Verification of a Realistic Compiler.}
Communications of the ACM, 52(7):107--115, 2009.
\href{https://doi.org/10.1145/1538788.1538814}{doi:10.1145/1538788.1538814}.

\bibitem{ouroboros25}
S.\ P.\ Lutar Jr.
\emph{Governed Post-Determinism: A Unified Theory of Verifiable Autonomy
(Ouroboros Thesis).}
SZL Holdings, 2026. Concept DOI
\url{https://doi.org/10.5281/zenodo.19944926}.
\end{thebibliography}

\end{document}
